    В современной математике евклидова геометрия возникает как результат последовательного обогащения абстрактного
    множества дополнительными математическими структурами:
    \begin{gather*}
        \text{\textbf{Метрическое пространство}} \;\longrightarrow\; \text{\textbf{Нормированное пространство}}
    \;\longrightarrow\;\\
        \text{\textbf{Предгильбертово пространство}} \;\longrightarrow\; \text{\textbf{Евклидово
    пространство}}\\
    \end{gather*}

    \subsection{Метрическое пространство}
    Множество $M$ с функцией расстояния (метрикой) $d \colon M \times M \to \mathbb{R}_{\ge 0}$, удовлетворяющей трём
    аксиомам:
    \begin{enumerate}
        \item Тождество неразличимых: $d(x, y) = 0 \iff x = y$;
        \item Симметрия: $d(x, y) = d(y, x)$;
        \item Неравенство треугольника: $d(x, z) \le d(x, y) + d(y, z)$.
    \end{enumerate}
    В метрическом пространстве нет понятий «направление», «ноль» или «сложение» — есть только расстояния между точками.

    \subsection{Нормированное векторное пространство}
    Векторное пространство $V$, на котором задана \textbf{норма} $\|\cdot\| \colon V \to \mathbb{R}_{\ge 0}$,
    удовлетворяющая аксиомам:
    \begin{enumerate}
        \item Положительная определенность: $\|x\| \ge 0$, причём $\|x\| = 0 \iff x = 0$;
        \item Абсолютная однородность: $\|\alpha x\| = |\alpha| \, \|x\|$ для всех $\alpha \in \mathbb{R}, x \in V$;
        \item Неравенство треугольника (субаддитивность): $\|x + y\| \le \|x\| + \|y\|$.
    \end{enumerate}
    \textbf{Индуцирование метрики:} Всякая норма естественным образом порождает метрику расстояния:
    \[ d(x, y) \coloneqq \|x - y\|. \]
    Обратное неверно: не всякая метрика порождается нормой (например, дискретная метрика $d(x, y) = \begin{cases}
                                                                                                        0, & x=y
                                                                                                        \\ 1, & x \ne y
    \end{cases}$ не удовлетворяет условию однородности).

    \subsection{Пространство со скалярным произведением (предгильбертово)}
    Векторное пространство $V$, оснащенное билинейной операцией $\langle \cdot, \cdot \rangle \colon V \times V \to
    \mathbb{R}$, удовлетворяющей аксиомам:
    \begin{enumerate}
        \item Симметричность: $\langle x, y \rangle = \langle y, x \rangle$;
        \item Линейность по первому аргументу: $\langle \alpha x + \beta y, z \rangle = \alpha \langle x, z \rangle +
        \beta \langle y, z \rangle$;
        \item Положительная определенность: $\langle x, x \rangle \ge 0$, причём $\langle x, x \rangle = 0 \iff x = 0$.
    \end{enumerate}

    \begin{theorem}[Неравенство Коши-Буняковского-Шварца]
Для любых элементов $x, y \in V$ со скалярным произведением справедливо:
\[ |\langle x, y \rangle| \le \sqrt{\langle x, x \rangle} \, \sqrt{\langle y, y \rangle}. \]
    \end{theorem}

    \textbf{Индуцирование нормы:} Всякое скалярное произведение порождает евклидову норму:
    \[ \|x\| \coloneqq \sqrt{\langle x, x \rangle}. \]
    Проверка неравенства треугольника для этой нормы опирается непосредственно на неравенство Коши-Буняковского-Шварца:
    \[ \|x + y\|^2 = \langle x+y, x+y \rangle = \|x\|^2 + 2\langle x, y \rangle + \|y\|^2 \le \|x\|^2 + 2\|x\|\|y\| +
    \|y\|^2 = (\|x\| + \|y\|)^2. \]

    \begin{theorem}[Критерий согласованности нормы со скалярным произведением: тождество параллелограмма]
Норма $\|\cdot\|$ на векторном пространстве $V$ порождается некоторым скалярным произведением тогда и только тогда,
когда для любых $x, y \in V$ выполняется \textbf{тождество параллелограмма}:
\[ \|x + y\|^2 + \|x - y\|^2 = 2\|x\|^2 + 2\|y\|^2. \]
Если тождество параллелограмма выполнено, скалярное произведение однозначно восстанавливается по поляризационному
тождеству:
\[ \langle x, y \rangle = \frac{1}{4}\left(\|x + y\|^2 - \|x - y\|^2\right). \]
    \end{theorem}

    \begin{example}[Контрпример: нормы, не порождаемые скалярным произведением]
Рассмотрим пространство $\mathbb{R}^2$:
\begin{itemize}
    \item Манхэттенская норма ($L_1$): $\|x\|_1 = |x_1| + |x_2|$.
    Возьмём $x = (1, 0)$ и $y = (0, 1)$.
    Тогда $\|x\|_1 = 1$, $\|y\|_1 = 1$.
    Имеем $x+y = (1, 1)$, $x-y = (1, -1)$, отсюда $\|x+y\|_1^2 + \|x-y\|_1^2 = 2^2 + 2^2 = 8$.
    Однако $2\|x\|_1^2 + 2\|y\|_1^2 = 2(1) + 2(1) = 4$.
    Так как $8 \ne 4$, норма $L_1$ не порождается скалярным произведением.
    \item Чебышёвская норма ($L_\infty$): $\|x\|_\infty = \max(|x_1|, |x_2|)$.
    Аналогично не удовлетворяет тождеству параллелограмма.
\end{itemize}
    \end{example}

    \subsection{Евклидово пространство}
    \textbf{Евклидовым пространством} называется конечномерное вещественное векторное пространство, оснащенное скалярным
    произведением.
    В нём одновременно согласованы все уровни геометрии:
    \[ \begin{array}{lcccc}
           \text{Скалярное произведение:} & \langle x, y \rangle & & & \\
           & \Downarrow & & & \\
           \text{Длины (норма):} & \|x\| = \sqrt{\langle x, x \rangle} & & & \\
           & \Downarrow & \searrow & & \\
           \text{Расстояния (метрика):} & d(x, y) = \|x - y\| & & \text{Углы:} & \cos \theta = \frac{\langle x, y
           \rangle}{\|x\| \|y\|}
    \end{array} \]